
\documentclass[addpoints]{exam}
\usepackage{amsmath}
\usepackage{amssymb}
\usepackage{color}
\usepackage{siunitx}

% \printanswers

\newcommand{\event}{Final Exam}
\newcommand{\tournament}{}
\def\type{0}

\setlength\fillinlinelength{\textwidth}
\CorrectChoiceEmphasis{\color{red}\bfseries}
\SolutionEmphasis{\color{red}\bfseries}

\setlength\answerclearance{2pt}
\pagestyle{headandfoot}
\runningheadrule
\runningheader{SPH3U \event}{\tournament}{\ifnum\type=1 Class Set Test \else Full Name:\hspace{1cm} \fi}
\runningfootrule
\runningfooter{}{Page \thepage\ of \numpages}{}

\begin{document}
\begin{center}
    {\huge Grade 11 Physics \\ \event
    \ifprintanswers \space Key
    \else \space \ifnum\type<2 \fi
    \fi \\ \vspace{3mm}\tournament\vspace{1mm}}
\end{center}

\vspace{.20\textheight}

\begin{center}
    First Name: \ifprintanswers
    \underline{\hspace{3.5cm}}\textbf{KEY}\underline{\hspace{5.5cm}}\\\vspace{5mm}
    \else \underline{\hspace{10cm}}\\ \vspace{5mm} \fi
    Last Name: \underline{\hspace{10cm}}\\ \vspace{5mm}
\end{center}

\noindent\textbf{\underline{Directions:}}
\begin{itemize}
    \item Show all work with units. FBDs required where applicable.
    \item $g = \SI{9.8}{m/s^2}$;\quad $v_\text{sound} = \SI{343}{m/s}$;\quad
    $k = \SI{8.99e9}{N\cdot m^2/C^2}$;\quad $e = \SI{1.60e-19}{C}$.
    \item This exam covers all units of SPH3U. Designed for 2 hours.
\end{itemize}
\begin{center}
\textit{For grading use only}\\
\multirowgradetable{2}[pages]
\end{center}

\newpage
\begin{questions}

\section*{Part A --- Multiple Choice (20 marks)}
\vspace{3mm}
\noindent\textit{Unit 1: Kinematics}
\vspace{3mm}

\question[1] A car slows from \SI{25}{m/s} to \SI{10}{m/s} in \SI{3.0}{s}. Its acceleration is:
\begin{choices}
    \choice \SI{5.0}{m/s^2} forward
    \correctchoice \SI{-5.0}{m/s^2} (deceleration)
    \choice \SI{8.3}{m/s^2}
    \choice \SI{11.7}{m/s^2}
\end{choices}

\question[1] A ball is thrown horizontally at \SI{10}{m/s} from a \SI{80}{m} cliff.
The time to land is:
\begin{choices}
    \choice \SI{2.0}{s}
    \correctchoice \SI{4.04}{s}
    \choice \SI{8.0}{s}
    \choice \SI{16}{s}
\end{choices}

\question[1] The range of the ball in the previous question is:
\begin{choices}
    \choice \SI{20}{m}
    \correctchoice \SI{40.4}{m}
    \choice \SI{80}{m}
    \choice \SI{160}{m}
\end{choices}

\question[1] Which graph shows uniform acceleration from rest?
\begin{choices}
    \choice A horizontal line on a $d$--$t$ graph
    \correctchoice A straight line with positive slope on a $v$--$t$ graph
    \choice A curved line on a $v$--$t$ graph
    \choice A horizontal line on a $v$--$t$ graph
\end{choices}

\vspace{3mm}
\noindent\textit{Unit 2: Forces}
\vspace{3mm}

\question[1] A \SI{8.0}{kg} box is pushed along a floor with a net force of \SI{20}{N}.
Its acceleration is:
\begin{choices}
    \choice \SI{0.4}{m/s^2}
    \correctchoice \SI{2.5}{m/s^2}
    \choice \SI{160}{m/s^2}
    \choice \SI{12}{m/s^2}
\end{choices}

\question[1] A block on a frictionless surface is pulled by a \SI{30}{N} force at
\ang{30} above horizontal. The horizontal component of this force is:
\begin{choices}
    \choice \SI{15}{N}
    \correctchoice \SI{26}{N}
    \choice \SI{30}{N}
    \choice \SI{17.3}{N}
\end{choices}

\question[1] A spring extends \SI{8.0}{cm} when a \SI{4.0}{N} force is applied.
Its spring constant is:
\begin{choices}
    \choice \SI{0.50}{N/m}
    \choice \SI{32}{N/m}
    \correctchoice \SI{50}{N/m}
    \choice \SI{200}{N/m}
\end{choices}

\question[1] Newton's Third Law: a book rests on a table. The reaction force to the
book's weight is:
\begin{choices}
    \choice The normal force from the table on the book
    \correctchoice The gravitational force the book exerts on Earth
    \choice The friction between table and floor
    \choice The normal force the book exerts on the table
\end{choices}

\vspace{3mm}
\noindent\textit{Unit 3: Energy}
\vspace{3mm}

\question[1] A \SI{2.0}{kg} object falls \SI{10}{m} from rest (no air resistance).
Its speed just before impact is:
\begin{choices}
    \choice \SI{10}{m/s}
    \correctchoice \SI{14}{m/s}
    \choice \SI{20}{m/s}
    \choice \SI{196}{m/s}
\end{choices}

\question[1] A force of \SI{40}{N} acts at \ang{90} to displacement. The work done is:
\begin{choices}
    \choice \SI{40}{J}
    \correctchoice \SI{0}{J}
    \choice \SI{-40}{J}
    \choice Depends on displacement
\end{choices}

\question[1] An engine does \SI{12\,000}{J} of work in \SI{60}{s}. Its power output is:
\begin{choices}
    \choice \SI{12\,000}{W}
    \correctchoice \SI{200}{W}
    \choice \SI{720\,000}{W}
    \choice \SI{20}{W}
\end{choices}

\question[1] A machine has an efficiency of 75\%. If the useful output is \SI{300}{J},
the energy input was:
\begin{choices}
    \choice \SI{225}{J}
    \correctchoice \SI{400}{J}
    \choice \SI{375}{J}
    \choice \SI{300}{J}
\end{choices}

\vspace{3mm}
\noindent\textit{Unit 4: Waves \& Sound}
\vspace{3mm}

\question[1] A closed air column resonates at \SI{85}{Hz} (fundamental). Its first
overtone is at:
\begin{choices}
    \choice \SI{170}{Hz}
    \correctchoice \SI{255}{Hz}
    \choice \SI{340}{Hz}
    \choice \SI{85}{Hz}
\end{choices}

\question[1] A wave travels at \SI{6.0}{m/s} with a period of \SI{0.25}{s}.
Its wavelength is:
\begin{choices}
    \choice \SI{0.25}{m}
    \correctchoice \SI{1.5}{m}
    \choice \SI{24}{m}
    \choice \SI{6.0}{m}
\end{choices}

\vspace{3mm}
\noindent\textit{Unit 5: Electricity}
\vspace{3mm}

\question[1] Three \SI{9}{\ohm} resistors are connected in series to a \SI{27}{V} source.
The current through the circuit is:
\begin{choices}
    \choice \SI{9.0}{A}
    \correctchoice \SI{1.0}{A}
    \choice \SI{3.0}{A}
    \choice \SI{0.33}{A}
\end{choices}

\question[1] The equivalent resistance of \SI{4}{\ohm} and \SI{12}{\ohm} in parallel is:
\begin{choices}
    \choice \SI{16}{\ohm}
    \choice \SI{8}{\ohm}
    \correctchoice \SI{3}{\ohm}
    \choice \SI{6}{\ohm}
\end{choices}

\question[1] Doubling the voltage across a resistor (constant $R$):
\begin{choices}
    \choice Doubles the current but halves the power
    \correctchoice Doubles the current and quadruples the power
    \choice Quadruples the current and doubles the power
    \choice Halves the current
\end{choices}

\question[1] According to Coulomb's Law, if both charges are doubled and the distance
is also doubled, the force:
\begin{choices}
    \choice Stays the same
    \correctchoice Stays the same
    \choice Doubles
    \choice Quadruples
\end{choices}

\question[1] A proton moves parallel to a magnetic field. The magnetic force on it is:
\begin{choices}
    \choice Maximum
    \choice Perpendicular to its velocity
    \correctchoice Zero
    \choice In the direction of motion
\end{choices}

\question[1] The energy stored in a spring compressed \SI{0.10}{m} with $k = \SI{800}{N/m}$ is:
\begin{choices}
    \choice \SI{80}{J}
    \choice \SI{8.0}{J}
    \correctchoice \SI{4.0}{J}
    \choice \SI{40}{J}
\end{choices}


\section*{Part B --- Short Answer (40 marks)}

\question \textbf{(Kinematics)} A stone is thrown from a cliff \SI{60}{m} above the
ocean at \SI{15}{m/s} horizontally.
\begin{parts}
    \part[2] Find the time to reach the water.
    \part[2] Find the horizontal distance from the base of the cliff.
    \part[3] Find the speed and direction of the stone just before it hits the water.
\end{parts}
\begin{solution}[9cm]
\textbf{(a)} $\Delta y = \tfrac{1}{2}gt^2$:\quad
$t = \sqrt{\dfrac{2(60)}{9.8}} = \sqrt{12.24} = \mathbf{\SI{3.50}{s}}$

\textbf{(b)} $\Delta x = v_x t = 15 \times 3.50 = \mathbf{\SI{52.5}{m}}$

\textbf{(c)} $v_y = gt = 9.8 \times 3.50 = \SI{34.3}{m/s}$ down;\quad $v_x = \SI{15}{m/s}$
\[
v = \sqrt{15^2 + 34.3^2} = \sqrt{225 + 1176.5} = \sqrt{1401.5} \approx \mathbf{\SI{37.4}{m/s}}
\]
\[
\theta = \arctan\!\left(\frac{34.3}{15}\right) \approx \mathbf{\ang{66} \text{ below horizontal}}
\]
\end{solution}

\question \textbf{(Forces)} A \SI{15}{kg} box sits on a ramp inclined at \ang{25}.
The coefficient of static friction is $\mu_s = 0.45$.
\begin{parts}
    \part[2] Calculate the component of gravity parallel to the ramp (tending to slide
    the box down).
    \part[2] Calculate the maximum static friction force.
    \part[2] Will the box slide? Justify quantitatively.
    \part[2] If the box is given a push and begins to slide ($\mu_k = 0.30$), find its
    acceleration down the ramp.
\end{parts}
\begin{solution}[9cm]
\textbf{(a)} $F_\parallel = mg\sin\ang{25} = 15(9.8)(0.4226) = \mathbf{\SI{62.1}{N}}$

\textbf{(b)} $F_N = mg\cos\ang{25} = 15(9.8)(0.9063) = \SI{133.2}{N}$;\quad
$f_{s,\max} = \mu_s F_N = 0.45 \times 133.2 = \mathbf{\SI{59.9}{N}}$

\textbf{(c)} $F_\parallel = \SI{62.1}{N} > f_{s,\max} = \SI{59.9}{N}$. The gravitational
component exceeds maximum static friction, so the box \textbf{will slide}.

\textbf{(d)} $f_k = 0.30 \times 133.2 = \SI{39.96}{N}$;\quad
$F_\text{net} = 62.1 - 39.96 = \SI{22.1}{N}$ down slope
\[
a = \frac{22.1}{15} = \mathbf{\SI{1.48}{m/s^2} \text{ down the ramp}}
\]
\end{solution}

\question \textbf{(Energy)} A \SI{500}{kg} roller coaster car starts from rest at the
top of a \SI{40}{m} hill (point A). It rolls down to point B at the bottom
(\SI{0}{m}), then up to point C at height \SI{25}{m}.
Ignore friction.
\begin{parts}
    \part[2] Find the speed of the car at point B.
    \part[2] Find the speed of the car at point C.
    \part[3] In reality, the car reaches point B at only \SI{24}{m/s} due to friction.
    Find the energy lost to friction between A and B, and the average friction force
    over the \SI{60}{m} slope from A to B.
\end{parts}
\begin{solution}[9cm]
\textbf{(a)} $mgh_A = \tfrac{1}{2}mv_B^2$:\quad
$v_B = \sqrt{2gh_A} = \sqrt{2(9.8)(40)} = \sqrt{784} = \mathbf{\SI{28.0}{m/s}}$

\textbf{(b)} $\tfrac{1}{2}mv_B^2 = \tfrac{1}{2}mv_C^2 + mgh_C$:\quad
$v_C = \sqrt{v_B^2 - 2gh_C} = \sqrt{784 - 2(9.8)(25)} = \sqrt{784-490} = \sqrt{294}
= \mathbf{\SI{17.1}{m/s}}$

\textbf{(c)} Actual $E_k$ at B: $\tfrac{1}{2}(500)(24)^2 = \SI{144\,000}{J}$

Initial PE: $mgh_A = 500(9.8)(40) = \SI{196\,000}{J}$

Energy lost: $196\,000 - 144\,000 = \mathbf{\SI{52\,000}{J}}$

$W_\text{friction} = f_k \cdot d \Rightarrow f_k = \dfrac{52\,000}{60} = \mathbf{\SI{867}{N}}$
\end{solution}

\question \textbf{(Waves)} An open air column has a fundamental frequency of
\SI{220}{Hz} at \SI{20}{\celsius} (speed of sound = \SI{343}{m/s}).
\begin{parts}
    \part[2] Find the length of the air column.
    \part[2] Find the frequency of the second harmonic (first overtone).
    \part[3] A source moving toward this pipe at \SI{25}{m/s} emits sound at \SI{220}{Hz}.
    What frequency does a stationary listener near the pipe hear?
\end{parts}
\begin{solution}[8cm]
\textbf{(a)} Open column: $f_1 = \dfrac{v}{2L} \Rightarrow L = \dfrac{v}{2f_1}
= \dfrac{343}{2 \times 220} = \dfrac{343}{440} = \mathbf{\SI{0.780}{m}}$

\textbf{(b)} Open columns support all harmonics; second harmonic $= 2f_1 = \mathbf{\SI{440}{Hz}}$

\textbf{(c)} Source approaching, observer stationary:
\[
f' = f_s \cdot \frac{v}{v - v_s} = 220 \times \frac{343}{343-25}
= 220 \times \frac{343}{318} = \mathbf{\SI{237}{Hz}}
\]
\end{solution}

\question \textbf{(Electricity)} A circuit has a \SI{24}{V} battery connected to $R_1 =
\SI{6.0}{\ohm}$ in series with a parallel combination of $R_2 = \SI{4.0}{\ohm}$ and
$R_3 = \SI{12}{\ohm}$.
\begin{parts}
    \part[2] Find the equivalent resistance of $R_2$ and $R_3$ in parallel.
    \part[2] Find the total circuit resistance and total current from the battery.
    \part[2] Find the voltage drop across $R_1$ and across the parallel combination.
    \part[2] Find the current through $R_2$ and $R_3$.
\end{parts}
\begin{solution}[9cm]
\textbf{(a)} $\dfrac{1}{R_{23}} = \dfrac{1}{4.0} + \dfrac{1}{12} = \dfrac{3}{12} + \dfrac{1}{12}
= \dfrac{4}{12} = \dfrac{1}{3}$;\quad $R_{23} = \mathbf{\SI{3.0}{\ohm}}$

\textbf{(b)} $R_T = R_1 + R_{23} = 6.0 + 3.0 = \SI{9.0}{\ohm}$;\quad
$I_T = \dfrac{V}{R_T} = \dfrac{24}{9.0} = \mathbf{\SI{2.67}{A}}$

\textbf{(c)} $V_1 = I_T R_1 = 2.67 \times 6.0 = \mathbf{\SI{16.0}{V}}$;\quad
$V_{23} = 24 - 16.0 = \mathbf{\SI{8.0}{V}}$ (or $= I_T \times R_{23} = 2.67 \times 3.0 = \SI{8.0}{V}$)

\textbf{(d)} $I_2 = \dfrac{V_{23}}{R_2} = \dfrac{8.0}{4.0} = \mathbf{\SI{2.0}{A}}$;\quad
$I_3 = \dfrac{8.0}{12} = \mathbf{\SI{0.667}{A}}$

Check: $I_2 + I_3 = 2.0 + 0.667 = \SI{2.67}{A} = I_T$ \checkmark
\end{solution}

\end{questions}
\end{document}
